% ==================================================================
% Chapter 5 — Implementation / Designing   (writer W5)
% ==================================================================
\chapter{Implementation}\label{ch:impl}

\section{Chapter overview}

This chapter explains how the artefact specified in \cref{ch:srs} was designed and built. It has four components: the Drone Trajectory eXtractor (\dtx{}), which turns the drone videos into trajectories, waiting times and calibration data; a Simulation of Urban MObility (\sumo{}) model of the junction; the Virtual LiDAR Sensor (\vls{}), which imitates an imperfect pole-mounted Light Detection and Ranging (LiDAR) unit; and the Pressure-based Real-time Intersection Signal Management (\prism{}) controller, the main contribution.
\Cref{sec:impl-dtx,sec:impl-sumo,sec:impl-vls,sec:impl-prism} present each algorithm and \cref{sec:impl-baselines} the baselines. \Cref{sec:impl-tech} justifies the technology stack and \cref{sec:impl-key} presents the key code.

The \sumo{} network, the \vls{} and \prism{} were first developed and tested with \emph{mock} demand split evenly across the four approaches; the drone-derived demand was introduced last, replacing the mock routes and calibrating the model. The control law itself is purely mathematical; artificial intelligence is used only for
vehicle detection inside \dtx{}.

% ------------------------------------------------------------------
\section{Design of the algorithms}\label{sec:impl-design}

\subsection{Drone Trajectory eXtractor (DTX)}\label{sec:impl-dtx}

\dtx{} converts each one-hour 4K video (30 frames per second, processed at 10) into metric trajectories and calibration quantities (\cref{alg:dtx}).

\paragraph{Stabilisation.} Wind caused slow drift of the hovering drone, so every frame was registered to a reference frame with Oriented FAST and Rotated BRIEF
(ORB) features~\cite{rublee2011orb}, filtered by Random Sample
Consensus (RANSAC)~\cite{fischler1981ransac}.   The estimated $3\times3$ matrix
$\mathbf{H}_t$ warps frame $t$ onto the reference.

\paragraph{Detection.} Vehicles were detected with YOLOv8s~\cite{jocher2023yolov8}, a single-stage
You Only Look Once (YOLO) detector, initialised from weights trained on the aerial VisDrone2019-DET
dataset~\cite{zhu2022visdrone}, whose \emph{tricycle} and \emph{awning-tricycle} classes resemble Sri
Lankan three-wheelers. The model was fine-tuned on \est{450} annotated frames from the two sessions
with the seven classes of \cref{tab:impl-vtypes}. At about \est{90~m} altitude a motorcycle occupies
roughly 20 pixels, so inference used Slicing Aided Hyper Inference (SAHI)~\cite{akyon2022sahi} with
$2\times2$ tiles and an overlap ratio of 0.2.

\paragraph{Tracking.} Detections were linked into tracks with ByteTrack~\cite{zhang2022bytetrack}
(track threshold 0.5, match threshold 0.8, buffer 30 frames).  Its second pass over low-confidence boxes recovers partly hidden motorcycles.
The class of a track was fixed by majority vote over its frames.

\paragraph{Pixel-to-metre conversion.} A ground-plane homography was estimated from at least four
control points (stop-line ends and lane-marking corners, using the 3.5~m lane width and satellite
imagery). A stabilised pixel $(u,v)$ maps to ground coordinates $(x,y)$ by
\begin{equation}
  \lambda \begin{bmatrix} x \\ y \\ 1 \end{bmatrix} = \mathbf{G}
  \begin{bmatrix} u \\ v \\ 1 \end{bmatrix}, \qquad \mathbf{G}\in\mathbb{R}^{3\times3},
  \label{eq:impl-homography}
\end{equation}
where $\lambda$ is a scale factor~\cite{hartley2004mvg}. The resulting ground sample distance (GSD)
was \est{$\approx$2.6~cm/px} and the reprojection error on held-out points \est{0.18~m}.

\paragraph{Smoothing and speed.} Positions were smoothed with a Savitzky--Golay filter
\cite{savitzky1964} (window of 11 samples, polynomial order 2), which removes jitter without the lag of a moving average. The speed at sample $k$ is
\begin{equation}
  v_k = \frac{\lVert \tilde{\mathbf{p}}_k - \tilde{\mathbf{p}}_{k-1} \rVert}{\Delta t},
  \qquad \Delta t = 0.1~\text{s},
  \label{eq:impl-speed}
\end{equation}
where $\tilde{\mathbf{p}}_k$ is the smoothed position. 

\paragraph{Zoning and movements.} Polygons were drawn for each approach, each exit and each slip
lane, together with stop-line segments. The entry and exit polygons crossed by a track define its
turning movement (left via slip lane, through, right or U-turn).

\paragraph{Waiting time and delay.} As in \sumo{}, a vehicle is waiting in any sample where its speed is
below 0.1~m/s while it is inside an approach zone $\mathcal{Z}$. The waiting time of vehicle $i$ is
\begin{equation}
  W_i = \sum_{k} \Delta t \cdot \mathbb{1}\!\left[ v_k < 0.1~\text{m/s} \;\wedge\; \tilde{\mathbf{p}}_k \in \mathcal{Z} \right],
  \label{eq:impl-wait}
\end{equation}
and its control delay is
\begin{equation}
  d_i = T_{\text{frame},i} - \frac{L_{\text{path},i}}{v_{\text{free}}},
  \label{eq:impl-delay}
\end{equation}
where $T_{\text{frame},i}$ is the time spent in the field of view, $L_{\text{path},i}$ the length of the
smoothed path and $v_{\text{free}}$ the observed free-flow speed of \est{38~km/h}. Tracks that were
already present when the video began, or still present when it ended, are \emph{censored}: they were counted in flows but excluded from averages.

\paragraph{Derived quantities.} Flows per approach and class were aggregated into 5-minute bins in
vehicles per hour and in Passenger Car Units (PCU) per hour,  with turning ratios and queue lengths. Saturation flow was estimated from queue-discharge headways
at the start of green, ignoring the first three vehicles
\cite{hcm2016}. Expressed per PCU,
\begin{equation}
  s = 3600 \cdot \frac{\sum_{j\geq4} \text{PCU}_j}{\sum_{j\geq4} h_j},
  \label{eq:impl-satflow}
\end{equation}
where $h_j$ is the stop-line headway of the $j$-th discharging vehicle. Pooled over both sessions this
gave \est{1850~PCU/h/lane}.

\paragraph{Output.} Finally, \dtx{} wrote a \sumo{} route file (a vehicle-type distribution per session
and one \texttt{<flow>} per movement per 5-minute bin) and a comma-separated values (CSV) file of
observed key performance indicators (KPIs) for calibration. In the morning the observed demand
was \est{3,420~veh/h}, or \est{3,150~PCU/h}, and in the afternoon it was \est{2,280~veh/h}, or
\est{2,085~PCU/h}. The observed average waiting times under the existing fixed-time plan were
\est{104.7~s} and \est{46.3~s} respectively.

\begin{algorithm}[htbp]
\DontPrintSemicolon
\KwIn{video $V$ (30 fps), zone polygons $\mathcal{Z}$, control points, model weights}
\KwOut{trajectory table $\mathcal{T}$, route file \texttt{.rou.xml}, observed-KPI CSV}
$F_0 \gets$ reference frame; $\mathcal{T} \gets \emptyset$\;
\ForEach{frame $F_t$ of $V$ sampled at 10 fps}{
  $\mathbf{H}_t \gets$ \textsc{RANSAC}(\textsc{ORBMatch}$(F_t, F_0)$); $\hat F_t \gets$ \textsc{Warp}$(F_t, \mathbf{H}_t)$\;
  $D_t \gets$ \textsc{SAHI}(YOLOv8s, $\hat F_t$, tiles $2\times2$, overlap 0.2)\;
  $\mathcal{T} \gets$ \textsc{ByteTrack}$(\mathcal{T}, D_t)$\;
}
\ForEach{track $\tau \in \mathcal{T}$}{
  class$(\tau) \gets$ majority vote of per-frame classes\;
  map pixels to metres with $\mathbf{G}$ (\cref{eq:impl-homography}); smooth (Savitzky--Golay 11/2)\;
  compute $v_k$ (\cref{eq:impl-speed}), $W_\tau$ (\cref{eq:impl-wait}) and $d_\tau$ (\cref{eq:impl-delay})\;
  movement$(\tau) \gets$ (entry zone, exit zone)\;
  \lIf{$\tau$ touches the video start or end}{mark $\tau$ censored}
}
aggregate flows (veh/h, PCU/h) per approach, class and 5-min bin; turning ratios; queues\;
$s \gets$ \cref{eq:impl-satflow} over discharge headways of 4th and later vehicles\;
write \texttt{vTypeDistribution} and \texttt{<flow>} elements; write KPI CSV (uncensored averages)\;
\caption{\dtx{} drone trajectory extraction}\label{alg:dtx}
\end{algorithm}

\subsection{SUMO network and demand}\label{sec:impl-sumo}

The network was exported from OpenStreetMap (OSM)~\cite{haklay2008osm} with \texttt{osmWebWizard},
converted with \texttt{netconvert --lefthand} because Sri Lanka drives on the left, and cleaned by hand
in \texttt{netedit}. Each arm has two entry and two exit lanes; the kerb lane carries through traffic
and the offside lane carries through and right-turning traffic. Each corner's channelised left-turn
slip lane was modelled as a separate edge that merges into the exit arm through an unsignalised
\emph{yield} link. The slip lanes are therefore not part of the Traffic Light System (TLS) program and
are ignored by the controller except as downstream exit occupancy. The central node is a
\texttt{traffic\_light} junction with the custom four-phase program of \cref{app:phases}: $\phi_1$
bypass through, $\phi_2$ bypass protected right, $\phi_3$ cross-road through and $\phi_4$ cross-road
protected right, each serving both opposing directions and each followed by 3~s amber and 2~s all-red.

Heterogeneous traffic was represented with the sublane model~\cite{sumo_sublane} at a lateral
resolution of 0.8~m. Motorcycles were given a high \texttt{lcSublane} eagerness and
\texttt{minGapLat} of 0.3~m so that they filter between queued cars to the stop line, which is the
dominant behaviour at Maviththara. Longitudinal motion used the Krauss model~\cite{krauss1998} with
driver imperfection $\sigma = 0.5$ and reaction time $\tau = 1.0$~s, and lateral decisions used the
model of Erdmann~\cite{erdmann2015lanechange}. A custom \texttt{threewheeler} type was created because
\sumo{} has no such class. The vehicle types are listed in \cref{tab:impl-vtypes}; the PCU values are
adopted values for urban mixed traffic~\cite{irc106}.

\begin{table}[htbp]
\centering
\caption{\sumo{} vehicle types used for Sri Lankan mixed traffic}\label{tab:impl-vtypes}
\small
\begin{tabular}{@{}llrrrrr@{}}
\toprule
Class & vClass & Len.\ (m) & PCU & Max (km/h) & Accel.\ (m/s$^2$) & Share M/A (\%)\\
\midrule
Motorcycle & \texttt{motorcycle} & 2.0 & 0.5 & 70 & 2.8 & \est{38} / \est{34}\\
Three-wheeler & \texttt{passenger}$^{*}$ & 2.7 & 0.8 & 50 & 1.8 & \est{24} / \est{27}\\
Car & \texttt{passenger} & 4.5 & 1.0 & 70 & 2.6 & \est{21} / \est{22}\\
Van/SUV & \texttt{delivery} & 5.0 & 1.0 & 70 & 2.6 & \est{7} / \est{8}\\
Bus & \texttt{bus} & 11.5 & 3.0 & 60 & 1.2 & \est{6} / \est{5}\\
Lorry & \texttt{truck} & 8.5 & 2.5 & 55 & 1.3 & \est{3} / \est{3}\\
Bicycle & \texttt{bicycle} & 1.8 & 0.3 & 20 & 1.2 & \est{1} / \est{1}\\
\bottomrule
\multicolumn{7}{@{}l}{\footnotesize $^{*}$Custom type \texttt{threewheeler} (width 1.4~m). M/A = morning/afternoon.}
\end{tabular}
\end{table}

Demand was produced in two stages. First, a mock-data generator wrote route files with equal demand
on every approach at 300, 500 and 700~veh/h, turning shares of 25\% left, 50\% through and 25\% right,
and a composition of 40\% motorcycles, 20\% three-wheelers, 25\% cars, 7\% vans, 5\% buses and 3\%
lorries. These were used for unit tests, phase-safety checks and
the parameter sweep. Second, the \dtx{} route files for the morning and afternoon sessions replaced the
mock routes. Simulations used a step length of 0.5~s, a controller decision every 1~s, a 300~s
warm-up that was discarded and 3,600~s of measurement, with tuning seeds 1--5 and evaluation seeds
6--15. Trip information, summary, queue and emission outputs (HBEFA4~\cite{hbefa41}) were recorded.
The calibration of \texttt{sigma}, \texttt{tau}, \texttt{minGap}, \texttt{lcSublane} and the speed
factor distribution against the observed counts and waiting times is reported in
\cref{sec:eval-calibration}.

\subsection{Virtual LiDAR Sensor (VLS)}\label{sec:impl-vls}

The \vls{} reproduces the output of a 360\textdegree{} unit,
comparable to a 32-channel Ouster OS1-32~\cite{ouster_os1}, mounted on a 6~m pole at the junction centre. Rather than simulating individual laser returns, \vls{} works at the \emph{object-list} level produced by roadside LiDAR perception software~\cite{zhao2019lidar,lv2019lidar}. Every second it queries all vehicles
through the Traffic Control Interface (TraCI)~\cite{wegener2008traci}, keeps those within
$R = 75$~m and applies the error model below (\cref{alg:vls}).

Position noise grows with range $r$:
\begin{equation}
  \hat{\mathbf{p}} = \mathbf{p} + \boldsymbol{\eta}, \qquad
  \boldsymbol{\eta} \sim \mathcal{N}\!\left(\mathbf{0}, \sigma(r)^2\mathbf{I}\right), \qquad
  \sigma(r) = 0.10 + 0.002\,r~\text{m}.
  \label{eq:impl-noise}
\end{equation}
Each object is independently missed with probability
\begin{equation}
  p_{\text{miss}}(r) = 0.01 + 0.04\left(\frac{r}{R}\right)^{2},
  \label{eq:impl-dropout}
\end{equation}
which reflects the thinning of beams with distance. Occlusion is computed by ray casting in the
horizontal plane. Let the interval $[\alpha_i^-, \alpha_i^+]$ be the azimuth range that object $i$
subtends at the sensor. Object $i$ is hidden if a nearer object $j$ that is at least 2.5~m tall
(a bus or lorry) satisfies
\begin{equation}
  \frac{\big|[\alpha_i^-, \alpha_i^+] \cap [\alpha_j^-, \alpha_j^+]\big|}{\alpha_i^+ - \alpha_i^-} \;\geq\; 0.7 .
  \label{eq:impl-occlusion}
\end{equation}
 Clutter is
added as $N_c \sim \text{Poisson}(\lambda = 0.2)$ false objects per scan at uniformly random positions
within range.

A real LiDAR does not know vehicle classes, so \vls{} reports only a \emph{size class} derived from the
measured length: \emph{small} below 3.2~m (motorcycles, bicycles and three-wheelers) mapped to
0.6~PCU, \emph{medium} from 3.2 to 6~m (cars and vans) mapped to 1.0~PCU, and \emph{large} above
6~m (buses and lorries) mapped to 2.8~PCU.  The small and large values are demand-weighted compromises between the classes they contain. The controller never sees the true class.

Raw detections then pass through a perception tracker that associates each detection with the nearest
predicted track within a 2.5~m gate and filters it with an $\alpha$--$\beta$ filter:
\begin{align}
  \hat{\mathbf{p}}^{-}_{k} &= \hat{\mathbf{p}}_{k-1} + \hat{\mathbf{v}}_{k-1}\,\Delta t, \label{eq:impl-ab1}\\
  \hat{\mathbf{p}}_{k} &= \hat{\mathbf{p}}^{-}_{k} + \alpha\,(\mathbf{z}_k - \hat{\mathbf{p}}^{-}_{k}), \label{eq:impl-ab2}\\
  \hat{\mathbf{v}}_{k} &= \hat{\mathbf{v}}_{k-1} + \frac{\beta}{\Delta t}\,(\mathbf{z}_k - \hat{\mathbf{p}}^{-}_{k}), \label{eq:impl-ab3}
\end{align}
with $\alpha = 0.5$, $\beta = 0.1$ and $\Delta t = 1$~s.  A track is confirmed after two
hits and deleted after three misses, which suppresses most single-scan clutter. Confirmed tracks are
map-matched to an approach and lane by point-in-polygon tests, and their distance to the stop line is
measured along the lane.  This list is the only input to \prism{}, so a real sensor could replace it unchanged.

\begin{algorithm}[htbp]
\DontPrintSemicolon
\KwIn{ground-truth vehicles $\mathcal{V}$ from TraCI, perception tracks $\mathcal{K}$}
\KwOut{object list $\mathcal{O}$ for the controller}
$\mathcal{D} \gets \{ i \in \mathcal{V} : r_i \leq R \}$\;
\ForEach{$i \in \mathcal{D}$ in increasing range}{
  \lIf{$\exists j : h_j \geq 2.5~\text{m},\, r_j < r_i$, overlap $\geq 0.7$ (\cref{eq:impl-occlusion})}{remove $i$; \textbf{continue}}
  \lIf{$U(0,1) < p_{\text{miss}}(r_i)$}{remove $i$; \textbf{continue}}
  $\hat{\mathbf{p}}_i \gets \mathbf{p}_i + \mathcal{N}(\mathbf{0}, \sigma(r_i)^2\mathbf{I})$; size class and PCU from length\;
}
add $N_c \sim \text{Poisson}(0.2)$ clutter objects uniformly within $R$\;
associate detections to $\mathcal{K}$ by nearest neighbour (gate 2.5~m); update with \cref{eq:impl-ab1,eq:impl-ab2,eq:impl-ab3}\;
confirm tracks after 2 hits; delete after 3 misses\;
$\mathcal{O} \gets$ confirmed tracks map-matched to approach, lane and $d_{\text{stop}}$\;
\caption{\vls{} scan and perception}\label{alg:vls}
\end{algorithm}

\subsection{PRISM control algorithm}\label{sec:impl-prism}

\prism{} belongs to the max-pressure family~\cite{varaiya2013maxpressure,tassiulas1992} but was adapted
to three realities of the Maviththara junction: vehicles differ greatly in size, a starved approach is
socially unacceptable even if it is theoretically optimal, and the sensor is imperfect.  It contains no learned component.

\subsubsection{State}

A \emph{movement} $m$ is a pair (approach, through or right). Left turns use the slip lanes and are
not controlled. With the lane assumption stated in \cref{sec:impl-sumo}, the controller computes for
each movement:
\begin{itemize}[noitemsep]
  \item $Q_m$, the queued PCU: the sum of PCU of objects in $m$'s lanes with $v < 1.4$~m/s within
        75~m of the stop line. Offside-lane objects are split between through and right in proportion
        to the \dtx{} turning shares, because the sensor cannot know a stationary vehicle's intention;
  \item $A_m$, the arriving PCU: the sum of PCU of moving objects whose estimated time of
        arrival (ETA) $d_{\text{stop}}/v$ is at most $H = 10$~s;
  \item $O_b$, the PCU occupancy of exit link $b$, whose capacity is $C_b \approx 23$~PCU
        ($2 \times 75~\text{m} / 6.5~\text{m}$ per PCU).
\end{itemize}
For each phase $\phi$, $R_\phi$ is the time since $\phi$ last had green.

\subsubsection{Equations}

The movement pressure is
\begin{equation}
  p_m = Q_m + \gamma A_m - \delta\, O_{\text{exit}(m)}, \qquad \gamma = 0.5,\; \delta = 0.6.
  \label{eq:pressure}
\end{equation}
The terms are the upstream queue, the platoon about to arrive and the downstream penalty that discourages feeding a congested exit
\cite{gregoire2015capacity}. Phase pressure aggregates its movements:
\begin{equation}
  P_\phi = \sum_{m \in \phi} p_m .
  \label{eq:phase-pressure}
\end{equation}
To encode fairness, each phase receives an urgency multiplier that grows with the square of its red
time,
\begin{equation}
  U_\phi = 1 + \kappa \left( \frac{R_\phi}{T_{\text{starve}}} \right)^{2}, \qquad \kappa = 1.5,
  \label{eq:urgency}
\end{equation}
so that urgency is negligible for short reds but reaches $1+\kappa = 2.5$ at the starvation cap. The
phase score is
\begin{equation}
  S_\phi = \max(P_\phi, 0)\cdot U_\phi .
  \label{eq:score}
\end{equation}
Clipping at zero stops an empty phase being promoted. The green time required to discharge the current demand is estimated as
\begin{equation}
  G_{\text{req}}(\phi) = \operatorname{clip}\!\left( l_s + h_s \cdot
  \frac{\sum_{m \in \phi} (Q_m + \gamma A_m)}{n_{\text{lanes}}(\phi)},\; G_{\min}(\phi),\; G_{\max}(\phi) \right),
  \label{eq:greq}
\end{equation}
with start-up lost time $l_s = 2.0$~s, saturation headway $h_s = 3600/1850 \approx 1.95$~s/PCU and
$n_{\text{lanes}}(\phi)$ the number of approach lanes discharging during $\phi$ (four for a through
phase, two for a right-turn phase).

\subsubsection{Decision loop}

Every $\Delta = 1$~s during the green of phase $g$ with elapsed green $t_g$, the controller evaluates
six rules in order and acts on the first that fires:
\begin{enumerate}[noitemsep]
  \item if $t_g < G_{\min}(g)$, continue;
  \item if $t_g \geq G_{\max}(g)$, terminate;
  \item if any other phase $\phi$ with waiting demand has $R_\phi \geq T_{\text{starve}}$, terminate
        (fairness override);
  \item \emph{gap-out}: if $\sum_{m\in g} Q_m \leq 0.5$~PCU and no object of $g$ has ETA
        $\leq \tau_{\text{gap}} = 3$~s, terminate;
  \item \emph{pressure switch}: if $\max_{\phi \neq g} S_\phi > (1+\theta) S_g + \varepsilon$ with
        $\theta = 0.25$, $\varepsilon = 2$~PCU, and $t_g \geq \rho\, G_{\text{req}}(g)$ with
        $\rho = 0.6$, terminate;
  \item otherwise extend the green by 1~s.
\end{enumerate}
On termination, 3~s amber and 2~s all-red follow. At the end of the all-red the next phase is chosen:
starved phases are served first, oldest first; otherwise the phase with the largest score $S_\phi$
($\phi \neq g$) is chosen, with ties broken in the order $\phi_1 \rightarrow \phi_2 \rightarrow \phi_3
\rightarrow \phi_4$.  Selection is therefore \emph{acyclic}. 

Two guards complete the design. The \emph{spillback guard} applies when any exit has
$O_b/C_b > 0.85$: every phase feeding $b$ has its effective maximum green reduced to $G_{\min}$, keeping vehicles out of a blocked junction box. The \emph{fail-safe} switches to the observed 78~s fixed-time plan
(B1) when the \vls{} delivers no data for five consecutive scans, and \prism{} resumes at the next
cycle boundary once ten consecutive valid scans have been received.  The complete logic is given in \cref{alg:prism} and
\cref{fig:impl-prism-flow}.

\begin{algorithm}[htbp]
\DontPrintSemicolon
\SetKwFunction{Decide}{Decide}\SetKwProg{Fn}{Function}{:}{}
\KwIn{object list $\mathcal{O}$, current green $g$, elapsed $t_g$, red timers $R_\phi$}
\KwOut{action $\in$ \{\textsc{Extend}, \textsc{Terminate}\} and a logged reason}
\Fn{\Decide{$\mathcal{O}, g, t_g$}}{
  \lIf{$\mathcal{O}$ missing for 5 consecutive scans}{\Return \textsc{FailSafe}(B1 plan)}
  compute $Q_m, A_m, O_b$ from $\mathcal{O}$; $p_m$ (\cref{eq:pressure}); $P_\phi$, $U_\phi$, $S_\phi$ (\cref{eq:phase-pressure,eq:urgency,eq:score})\;
  $G^{\text{eff}}_{\max} \gets G_{\min}(g)$ \textbf{if} $g$ feeds an exit with $O_b/C_b>0.85$ \textbf{else} $G_{\max}(g)$\;
  \lIf{$t_g < G_{\min}(g)$}{\Return (\textsc{Extend}, ``min green'')}
  \lIf{$t_g \geq G^{\text{eff}}_{\max}$}{\Return (\textsc{Terminate}, ``max out / spillback'')}
  \lIf{$\exists \phi \neq g: R_\phi \geq T_{\text{starve}}$ and $P_\phi > 0$}{\Return (\textsc{Terminate}, ``starvation'')}
  \lIf{$\sum_{m\in g} Q_m \leq 0.5$ and no ETA$_g \leq \tau_{\text{gap}}$}{\Return (\textsc{Terminate}, ``gap out'')}
  \If{$\max_{\phi\neq g} S_\phi > (1+\theta)S_g + \varepsilon$ \textbf{and} $t_g \geq \rho\,G_{\text{req}}(g)$}{
    \Return (\textsc{Terminate}, ``pressure switch'')}
  \Return (\textsc{Extend}, ``extend'')\;
}
\BlankLine
\textbf{On terminate:} amber 3 s; all-red 2 s; recompute $S_\phi$\;
\eIf{starved phases exist}{$g' \gets$ the starved phase with largest $R_\phi$}{$g' \gets \arg\max_{\phi\neq g} S_\phi$, ties $\phi_1\!\to\!\phi_2\!\to\!\phi_3\!\to\!\phi_4$}
start green of $g'$; $R_{g} \gets 0$ at end of $g$'s green; log decision record\;
\caption{\prism{} decision procedure (executed every $\Delta = 1$~s)}\label{alg:prism}
\end{algorithm}

\begin{figure}[htbp]
\centering
\resizebox{!}{0.66\textheight}{%
\begin{tikzpicture}[
  node distance=4.5mm and 9mm,
  box/.style={rectangle, rounded corners=2pt, draw, fill=black!5, align=center, font=\small, minimum width=38mm, minimum height=7mm},
  dec/.style={diamond, aspect=3.4, draw, align=center, font=\small, inner sep=0.5pt},
  term/.style={rectangle, rounded corners=7pt, draw, fill=black!12, align=center, font=\small, minimum width=26mm, minimum height=7mm},
  arr/.style={-{Stealth[length=2mm]}, thick},
  lab/.style={font=\footnotesize}]
\node[term] (start) {Every $\Delta=1$ s};
\node[box, below=of start] (scan) {Read \vls{} object list};
\node[dec, below=of scan] (fs) {Data missing\\for 5 scans?};
\node[box, below=of fs] (state) {Compute $Q_m, A_m, O_b$\\$p_m, P_\phi, U_\phi, S_\phi$};
\node[dec, below=of state] (gmin) {$t_g < G_{\min}$?};
\node[dec, below=of gmin] (gmax) {$t_g \geq G^{\text{eff}}_{\max}$?};
\node[dec, below=of gmax] (starve) {Other phase\\starved?};
\node[dec, below=of starve] (gap) {Gap-out?};
\node[dec, below=of gap] (press) {Pressure\\switch?};
\node[box, below=of press] (ext) {Extend green 1 s};
\coordinate (bus) at ($(gmax.east)+(22mm,0)$);
\node[box, right=14mm of fs] (fixed) {Fail-safe:\\B1 78 s plan};
\node[box, left=10mm of gmin] (cont) {Continue\\green};
\node[box] (term) at (bus |- ext) {Terminate $g$: amber\\3 s + all-red 2 s};
\node[box, below=of term] (sel) {Next: oldest starved,\\else $\arg\max S_\phi$\\(ties $\phi_1\!\to\!\phi_4$)};
\node[box, below=of ext] (log) {Log decision,\\reason and scores};
\draw[arr] (start) -- (scan);
\draw[arr] (scan) -- (fs);
\draw[arr] (fs) -- node[lab, right] {no} (state);
\draw[arr] (fs) -- node[lab, above] {yes} (fixed);
\draw[arr] (state) -- (gmin);
\draw[arr] (gmin) -- node[lab, right] {no} (gmax);
\draw[arr] (gmin) -- node[lab, above] {yes} (cont);
\draw[arr] (gmax) -- node[lab, right] {no} (starve);
\draw[arr] (starve) -- node[lab, right] {no} (gap);
\draw[arr] (gap) -- node[lab, right] {no} (press);
\draw[arr] (press) -- node[lab, right] {no} (ext);
\draw[thick] (gmax.east) -- node[lab, above] {yes} (bus);
\draw[thick] (starve.east) -- node[lab, above] {yes} (bus |- starve.east);
\draw[thick] (gap.east) -- node[lab, above] {yes} (bus |- gap.east);
\draw[thick] (press.east) -- node[lab, above] {yes} (bus |- press.east);
\draw[arr] (bus) -- (term.north);
\draw[arr] (term) -- (sel);
\draw[arr] (ext) -- (log);
\draw[arr] (sel.west) -- (sel.west -| log.east);
\draw[arr] (cont.south) |- (log.west);
\end{tikzpicture}}
\caption{Flowchart of the \prism{} decision procedure}\label{fig:impl-prism-flow}
\end{figure}

\subsubsection{Properties}

\paragraph{Safety.} \prism{} never writes a signal state directly. It selects one of four
pre-verified green states, and every transition passes through the 3~s amber and 2~s all-red
states of \cref{app:phases}, so conflicting greens and truncated intergreens are impossible by
construction. This was also checked by an automated test over every simulated step
(\cref{app:code}).

\paragraph{Bounded fairness.} Let phase $\phi$ end its green at time $t_0$. It becomes starved at
$t_0 + T_{\text{starve}}$. From then on, rule~3 terminates any running green at the first decision
step after that green reaches $G_{\min}$, and the oldest-first rule serves $\phi$ after the following
5~s intergreen. If another phase became starved earlier, it is served first, but its green is itself
cut at $G_{\min}$ because $\phi$ is starved. The extra time beyond $T_{\text{starve}}$ therefore
consists of at most minimum greens and intergreens, which together are always shorter than one
maximum through green plus one intergreen. A conservative bound on the red time of any phase is
\begin{equation}
  R^{\text{worst}} \leq T_{\text{starve}} + G_{\max} + 5~\text{s} = 120 + 50 + 5 = 175~\text{s},
  \label{eq:impl-fairness}
\end{equation}
using the through-phase $G_{\max}$ of 50~s, which is the longest green any phase can hold. 

\paragraph{Stability intuition.} Varaiya~\cite{varaiya2013maxpressure} proved that serving the phase
of maximum pressure stabilises queues for any demand inside the capacity region, using only local
information. \prism{} keeps that core because phases are ranked by pressure. However, the minimum and
maximum greens, the hysteresis $\theta$ and the urgency factor mean that the formal proof does not
transfer exactly.  This trade-off for fairness is also made by cycle-based variants~\cite{kouvelas2014maxpressure}.

\paragraph{Complexity.} Each object is assigned to one movement and one exit, so state computation is
$O(n)$ in the number of objects $n$.  The mean decision
time measured was \est{0.4~ms} (\cref{sec:eval-nonfunctional}).

\paragraph{Explainability.} Every decision writes a record containing time, phase, $t_g$, the rule
that fired, all $P_\phi$, $U_\phi$ and $S_\phi$, and $G_{\text{req}}$.  For example, ``pressure switch: $S_{\phi_3} = 18.4 > 1.25 \times
11.2 + 2$''. 

\paragraph{Parameter justification.} \Cref{tab:impl-prism-params} lists the design parameters and the
reasoning behind each value.  They were fixed on mock data and tuning seeds 1--5.

\begin{table}[htbp]
\centering
\caption{\prism{} design parameters and their justification}\label{tab:impl-prism-params}
\small
\begin{tabularx}{\textwidth}{@{}ll>{\raggedright\arraybackslash}X@{}}
\toprule
Parameter & Value & Justification\\
\midrule
$\gamma$ (arrival weight) & 0.5 & Arrivals within $H$ are uncertain and not yet delaying anyone, so they count half as much as queued PCU.\\
$\delta$ (exit weight) & 0.6 & Penalises congested exits without letting downstream noise dominate the upstream queue ($\delta<1$).\\
$\kappa$ (urgency gain) & 1.5 & Score at most $2.5\times$ at $T_{\text{starve}}$; the quadratic form keeps short reds almost unaffected.\\
$\theta$ (hysteresis) & 0.25 & Prevents oscillation from sensor noise; a challenger must be 25\% stronger.\\
$\varepsilon$ (margin) & 2 PCU & Absolute margin for light traffic, where relative margins are meaningless.\\
$\rho$ (discharge ratio) & 0.6 & Ensures most of the existing queue has discharged before a pressure switch.\\
$H$ (horizon) & 10 s & About 100 m at free-flow speed, which covers the 75 m sensor range.\\
$\tau_{\text{gap}}$ (gap time) & 3 s & Standard actuated gap; equal to the B2 \texttt{max-gap} for a fair comparison.\\
$l_s$ (start-up lost time) & 2.0 s & Typical start-up lost time~\cite{hcm2016}.\\
$h_s$ (headway per PCU) & 1.95 s & $3600/1850$, from the \dtx{} saturation flow.\\
$G_{\min}$ & 10 / 5 s & Through / right; driver expectation and clearance of the queue head.\\
$G_{\max}$ & 50 / 25 s & Through / right; limits red time of the other three phases.\\
$T_{\text{starve}}$ & 120 s & Longer than any observed red (67 s) and acceptable to waiting drivers.\\
Amber / all-red & 3 s / 2 s & Same as the existing signal.\\
$\Delta$ (decision period) & 1 s & Two \sumo{} steps; matches the \vls{} scan rate.\\
\bottomrule
\end{tabularx}
\end{table}

\subsection{Baseline controllers}\label{sec:impl-baselines}

Three baselines and an oracle variant were implemented.

\paragraph{B1: observed fixed-time plan.} The existing signal runs the 78~s plan of
\cref{tab:impl-b1}. The greens of $\phi_1$ and $\phi_2$ and the red of $\phi_3$ were timed from the
footage. Since every cycle contains four intergreens of $3 + 2 = 5$~s, the $\phi_3$ green is
$78 - 3 - 60 = 15$~s, and the remaining green follows from
$21 + 8 + 15 + G_{\phi_4} + 4 \times 5 = 78$, giving \est{14~s} for $\phi_4$.  

\begin{table}[htbp]
\centering
\caption{Baseline B1: observed fixed-time plan (cycle 78 s)}\label{tab:impl-b1}
\small
\begin{tabular}{@{}llrl@{}}
\toprule
Phase & Movement (both directions) & Green (s) & Source\\
\midrule
$\phi_1$ & Bypass through & 21 & Observed (red 54 s)\\
$\phi_2$ & Bypass protected right & 8 & Observed (red 67 s)\\
$\phi_3$ & Cross-road through & 15 & $78 - 3 - 60$\\
$\phi_4$ & Cross-road protected right & \est{14} & $78 - (21+8+15) - 4\times5$\\
\midrule
\multicolumn{2}{@{}l}{Intergreen after each phase} & 3 + 2 & Amber + all-red\\
\bottomrule
\end{tabular}
\end{table}

\paragraph{B2: SUMO actuated.} \sumo{}'s built-in gap-based actuated controller~\cite{sumo_actuated}
was run on the same four phases with \texttt{minDur} 10/5~s and \texttt{maxDur} 50/25~s for
through/right phases, the default induction loops and a maximum gap of 3~s.  It extends a green while vehicles keep arriving but ignores vehicle size and the demand on other phases.

\paragraph{B3: Webster-optimised fixed time.}  This separates adaptivity from simple re-timing. Webster's
method~\cite{webster1958} was applied to the \dtx{} hourly demand of each session (\cref{tab:impl-webster-demand}).
Left turns use the slip lanes and were excluded. Because waiting right-turners occupy the offside lane
during through phases, \dtx{} showed that \est{75\%} of through traffic used the kerb lane, so the
critical lane flow of a through phase is $0.75$ of the heavier approach's through flow. For a
right-turn phase it is the full right-turn flow in the offside lane. The flow ratio of phase $i$ is
$y_i = q_i/s$ with $s = $ \est{1850~PCU/h/lane}, $Y = \sum_i y_i$, the lost time is
$L = 4 \times 5 = 20$~s, and
\begin{equation}
  C_0 = \frac{1.5L + 5}{1 - Y}, \qquad g_i = (C_0 - L)\,\frac{y_i}{Y}.
  \label{eq:impl-webster}
\end{equation}

\begin{table}[htbp]
\centering
\caption{\dtx{} hourly demand used for B3 (PCU/h; L = slip lane, uncontrolled)}\label{tab:impl-webster-demand}
\small
\begin{tabular}{@{}lrrrrcrrrr@{}}
\toprule
 & \multicolumn{4}{c}{Morning} & & \multicolumn{4}{c}{Afternoon}\\
\cmidrule(lr){2-5}\cmidrule(lr){7-10}
Approach & L & T & R & Total & & L & T & R & Total\\
\midrule
Bypass 1 & \est{150} & \est{960} & \est{170} & \est{1,280} & & \est{90} & \est{570} & \est{120} & \est{780}\\
Bypass 2 & \est{120} & \est{680} & \est{130} & \est{930} & & \est{80} & \est{480} & \est{100} & \est{660}\\
Cross 1 & \est{100} & \est{340} & \est{150} & \est{590} & & \est{60} & \est{300} & \est{110} & \est{470}\\
Cross 2 & \est{60} & \est{210} & \est{80} & \est{350} & & \est{35} & \est{100} & \est{40} & \est{175}\\
\midrule
Total & \est{430} & \est{2,190} & \est{530} & \est{3,150} & & \est{265} & \est{1,450} & \est{370} & \est{2,085}\\
\bottomrule
\end{tabular}
\end{table}

\begin{table}[htbp]
\centering
\caption{Webster calculation for B3 (critical lane flows in PCU/h/lane, times in s)}\label{tab:impl-webster}
\small
\begin{tabular}{@{}llrrrclrrr@{}}
\toprule
 & & \multicolumn{3}{c}{Morning} & & & \multicolumn{3}{c}{Afternoon}\\
\cmidrule(lr){3-5}\cmidrule(lr){8-10}
Phase & Critical & $q_i$ & $y_i$ & $g_i$ & & Critical & $q_i$ & $y_i$ & $g_i$\\
\midrule
$\phi_1$ & Bypass 1 T & \est{720} & \est{0.389} & \est{54} & & Bypass 1 T & \est{428} & \est{0.231} & \est{23}\\
$\phi_2$ & Bypass 1 R & \est{170} & \est{0.092} & \est{13} & & Bypass 1 R & \est{120} & \est{0.065} & \est{6}\\
$\phi_3$ & Cross 1 T & \est{255} & \est{0.138} & \est{19} & & Cross 1 T & \est{225} & \est{0.122} & \est{12}\\
$\phi_4$ & Cross 1 R & \est{150} & \est{0.081} & \est{11} & & Cross 1 R & \est{110} & \est{0.059} & \est{6}\\
\midrule
\multicolumn{2}{@{}l}{$Y$} & & \est{0.700} & & & & & \est{0.477} & \\
\multicolumn{2}{@{}l}{$C_0$ (computed / used)} & & \multicolumn{2}{r}{\est{116.7} / \est{117}} & & & & \multicolumn{2}{r}{\est{66.9} / \est{67}}\\
\bottomrule
\end{tabular}
\end{table}

For the morning, $C_0 = 35/(1-0.700) = 116.7$~s, rounded to \est{117~s}. The effective green of
$117 - 20 = 97$~s splits as 53.90, 12.75, 19.12 and 11.22~s, implemented as \est{54/13/19/11~s}. For
the afternoon, $C_0 = 35/(1-0.477) = 66.9$~s, rounded to \est{67~s}. The 47~s of green splits as
22.76, 6.40, 12.02 and 5.81~s, implemented as \est{23/6/12/6~s}. Both plans respect the minimum
greens. The calculation explains the morning congestion: under B1, $\phi_1$ has a degree of
saturation $x = y\,C/g = 0.389 \times 78/21 \approx$ \est{1.44}, so the bypass through movement is
heavily oversaturated, while $\phi_4$ runs at only \est{0.45}. Webster also lengthens the morning
cycle to \est{117~s} and shortens the afternoon cycle to \est{67~s}, whereas the existing signal uses
78~s all day. The existing plan therefore allocates its green time to the wrong phases and uses one cycle length for very different demand levels.
However, Webster's formula assumes stationary random arrivals and undersaturation, and it ignores
queues that spill back past the slip-lane entries. The simulated benefit of B3 was therefore expected
to be smaller than the analytical one (\cref{sec:eval-results}).

\paragraph{PRISM-GT (oracle).} The same \prism{} code was run with the \vls{} replaced by an exact
ground-truth object list, with true positions, no dropout or occlusion and true-class PCU. The gap
between \prism{} and PRISM-GT measures the cost of imperfect sensing.

% ------------------------------------------------------------------
\section{Technology selection}\label{sec:impl-tech}

Every technology was selected by comparing realistic alternatives against the needs of this
research: free availability to a student, maturity, support for small objects in aerial video,
support for heterogeneous lane-less traffic and scriptable control of the signal.
\Cref{tab:impl-tech-core,tab:impl-tech-libs} summarise the decisions. No web front-end or custom
dashboard was built. Results and live decisions were presented through SUMO-GUI only, with the
\vls{} points and the current \prism{} reason overlaid.  This kept the effort on the control law.

\begin{table}[htbp]
\centering
\caption{Technology selection: language, models and simulator}\label{tab:impl-tech-core}
\small
\begin{tabularx}{\textwidth}{@{}>{\raggedright\arraybackslash}p{2.2cm}>{\raggedright\arraybackslash}p{3.1cm}>{\raggedright\arraybackslash}p{2.1cm}>{\raggedright\arraybackslash}X@{}}
\toprule
Concern & Alternatives & Chosen & Reason\\
\midrule
Programming language & C++, Java, MATLAB, Python & Python 3.11 & Native TraCI/libsumo bindings; Ultralytics, OpenCV and SciPy in one language; fast prototyping. C++ speed is unnecessary because a decision takes under 1 ms. MATLAB is licensed and has weak SUMO support.\\
Detection model & YOLO11, RT-DETR, Faster R-CNN, YOLOv8s & YOLOv8s & Mature, well documented, VisDrone weights available, trains at 1280 px on a free T4. YOLO11~\cite{jocher2024yolo11} offered marginal gains with less community evidence at the time; RT-DETR~\cite{zhao2024rtdetr} needs more GPU memory; Faster R-CNN~\cite{ren2015faster} is two-stage and too slow for two hours of video.\\
Tracker & SORT, DeepSORT, BoT-SORT, ByteTrack & ByteTrack & Recovers low-score boxes of small motorcycles; needs no re-identification network, which is of little use for top-down vehicles. SORT~\cite{bewley2016sort} loses occluded objects; DeepSORT~\cite{wojke2017deepsort} and BoT-SORT~\cite{aharon2022botsort} add appearance models with little benefit here.\\
Simulator & VISSIM, Aimsun, CARLA, SUMO & SUMO & Open source; sublane model for motorcycle filtering; TraCI for per-second signal control. VISSIM~\cite{fellendorf2010vissim} and Aimsun~\cite{casas2010aimsun} are commercial; CARLA~\cite{dosovitskiy2017carla} targets single-vehicle autonomy, not junction throughput.\\
\bottomrule
\end{tabularx}
\end{table}

\begin{table}[htbp]
\centering
\caption{Technology selection: libraries and tooling}\label{tab:impl-tech-libs}
\small
\begin{tabularx}{\textwidth}{@{}>{\raggedright\arraybackslash}p{2.6cm}>{\raggedright\arraybackslash}p{3.6cm}>{\raggedright\arraybackslash}X@{}}
\toprule
Need & Chosen (alternatives) & Reason\\
\midrule
Image processing & OpenCV~\cite{bradski2000opencv} (scikit-image) & ORB, RANSAC homography and warping in optimised C++ behind a Python interface.\\
Numerics & NumPy/SciPy~\cite{harris2020numpy,virtanen2020scipy} & Vectorised geometry; Savitzky--Golay filter; Wilcoxon test.\\
Tabular data & Pandas~\cite{mckinney2010pandas} (plain CSV) & Grouping trajectories by track, class and 5-min bin.\\
Tracking utilities & supervision~\cite{roboflow_supervision} & Maintained ByteTrack implementation, zone polygons and annotators.\\
Small objects & SAHI~\cite{akyon2022sahi} (higher input size only) & Tiled inference raises recall for 20-pixel motorcycles without retraining at 4K.\\
Simulation API & libsumo for batches, TraCI for GUI runs & libsumo runs in-process and avoids socket overhead; TraCI is required to drive SUMO-GUI.\\
Annotation & CVAT and Roboflow & Free; export directly to YOLO format; Roboflow handled the 70/15/15 split.\\
Testing & pytest (unittest) & Concise fixtures and parametrised tests for phase safety and edge cases.\\
Version control & Git and GitHub & History of every configuration and seed, which supports reproducibility.\\
Compute & Google Colab T4 GPU (local CPU) & Free GPU for fine-tuning; simulation runs on the student laptop.\\
Front-end & None (SUMO-GUI only) & A dashboard was a ``Won't have'' requirement; SUMO-GUI shows the network, vehicles and signal states.\\
\bottomrule
\end{tabularx}
\end{table}

% ------------------------------------------------------------------
\section{Key implementation}\label{sec:impl-key}

The code follows the repository layout planned in \cref{sec:srs-arch}: \texttt{extraction/} for \dtx{},
\texttt{sumo/} for the network and demand, \texttt{sensing/} for the \vls{}, \texttt{control/} for
\prism{} and the baselines, \texttt{experiments/} for batch runs and statistics, and \texttt{tests/}.
Only the fragments that embody the contribution are shown here; longer listings are in
\cref{app:code}.

\paragraph{Waiting-time computation.} \Cref{lst:impl-wait} implements \cref{eq:impl-speed,eq:impl-wait,eq:impl-delay}.
The decisive detail is that the 0.1~m/s threshold is the same constant \sumo{} uses, so observed and
simulated waiting times are directly comparable during calibration. 

\begin{lstlisting}[style=py,caption={Per-vehicle waiting time and control delay (\texttt{extraction/metrics.py})},label={lst:impl-wait}]
WAIT_SPEED = 0.1   # m/s, identical to SUMO's waitingTime threshold
DT = 0.1           # s, video processed at 10 fps

def vehicle_metrics(trk: pd.DataFrame, zones: ZoneMap, v_free: float) -> dict:
    """Waiting time W_i, time in frame and control delay d_i for one track."""
    xy = savgol_filter(trk[["x", "y"]].to_numpy(), 11, 2, axis=0)
    step = np.linalg.norm(np.diff(xy, axis=0), axis=1)
    v = np.r_[step[0], step] / DT                     # speed per sample
    in_approach = zones.in_approach(xy)               # point-in-polygon
    W = DT * np.count_nonzero((v < WAIT_SPEED) & in_approach)
    t_frame = trk.t.iloc[-1] - trk.t.iloc[0]
    delay = max(t_frame - step.sum() / v_free, 0.0)
    censored = (trk.t.iloc[0] <= zones.t_start + DT or
                trk.t.iloc[-1] >= zones.t_end - DT)
    return {"track_id": int(trk.track_id.iloc[0]),
            "cls": trk.cls.mode().iat[0],             # majority vote
            "movement": zones.movement(xy),
            "W": W, "t_frame": t_frame, "delay": delay,
            "censored": censored}
\end{lstlisting}

\paragraph{VLS scan.} \Cref{lst:impl-vls} shows the error model. Vehicles are processed in increasing
range so that the list of potential occluders only contains nearer tall vehicles, which keeps the
occlusion test cheap.

\begin{lstlisting}[style=py,caption={Noisy, occluded LiDAR scan (\texttt{sensing/virtual\_lidar.py})},label={lst:impl-vls}]
def scan(self, vehicles: list[GTVehicle]) -> list[Detection]:
    near = sorted((v for v in vehicles if self.range(v) <= self.R), key=self.range)
    dets, blockers = [], []                   # blockers: tall vehicles seen so far
    for v in near:
        r = self.range(v)
        span = self.azimuth_span(v)
        if any(overlap(span, b) >= 0.7 * width(span) for b in blockers):
            continue                          # occluded by bus/lorry
        if v.height >= 2.5:
            blockers.append(span)
        if self.rng.random() < 0.01 + 0.04 * (r / self.R) ** 2:
            continue                          # range-dependent dropout
        sigma = 0.10 + 0.002 * r
        x, y = v.pos + self.rng.normal(0.0, sigma, 2)
        length = v.length + self.rng.normal(0.0, sigma)
        dets.append(Detection(x, y, size_class(length)))
    for _ in range(self.rng.poisson(0.2)):    # clutter
        dets.append(Detection(*self.random_point_in_range(), "small"))
    return dets

def size_class(length: float) -> str:
    return "small" if length < 3.2 else "medium" if length <= 6.0 else "large"
\end{lstlisting}

\paragraph{Pressure and score.} \Cref{lst:impl-pressure} computes \cref{eq:pressure,eq:phase-pressure,eq:urgency,eq:score}
in a single pass over the object list. This is the source of the $O(n)$ complexity.

\begin{lstlisting}[style=py,caption={Movement pressure and phase score (\texttt{control/prism.py})},label={lst:impl-pressure}]
def movement_state(self, objs: list[TrackedObject]) -> dict[Movement, MovState]:
    st = {m: MovState() for m in MOVEMENTS}
    for o in objs:                                     # single O(n) pass
        if o.approach is None:                         # exit link
            self.occ[o.exit_link] += o.pcu
            continue
        for m, share in self.split(o):                 # offside lane: T/R by DTX shares
            if o.v < Q_SPEED and o.d_stop <= 75.0:
                st[m].Q += share * o.pcu
            elif o.v >= Q_SPEED and o.d_stop / o.v <= self.p.H:
                st[m].A += share * o.pcu
                st[m].min_eta = min(st[m].min_eta, o.d_stop / o.v)
    return st

def scores(self, st, t_now: float) -> dict[int, float]:
    S = {}
    for ph in PHASES:
        P = sum(st[m].Q + self.p.gamma * st[m].A
                - self.p.delta * self.occ[EXIT[m]] for m in ph.movements)
        U = 1.0 + self.p.kappa * (self.red_time(ph, t_now) / self.p.T_starve) ** 2
        S[ph.id] = max(P, 0.0) * U
    return S
\end{lstlisting}

\paragraph{Decision loop.} \Cref{lst:impl-decide} is a direct transcription of the six rules. Each branch returns a reason string, which provides the explainability.

\begin{lstlisting}[style=py,caption={The \texttt{decide()} rule cascade (\texttt{control/prism.py})},label={lst:impl-decide}]
def decide(self, objs, t_now: float) -> Decision:
    st = self.movement_state(objs)
    S = self.scores(st, t_now)
    g, tg, p = self.green, t_now - self.green_start, self.p
    gmax = g.gmin if self.spillback(g) else g.gmax
    if tg < g.gmin:
        return Decision(EXTEND, "min_green", S)
    if tg >= gmax:
        return Decision(TERMINATE, "max_out" if gmax == g.gmax else "spillback", S)
    if self.starved(st, t_now, exclude=g):
        return Decision(TERMINATE, "starvation", S)
    queue_g = sum(st[m].Q for m in g.movements)
    eta_g = min(st[m].min_eta for m in g.movements)
    if queue_g <= 0.5 and eta_g > p.tau_gap:
        return Decision(TERMINATE, "gap_out", S)
    best = max(S[k] for k in S if k != g.id)
    if best > (1 + p.theta) * S[g.id] + p.eps and tg >= p.rho * self.g_req(g, st):
        return Decision(TERMINATE, "pressure_switch", S)
    return Decision(EXTEND, "extend", S)
\end{lstlisting}

\paragraph{TraCI integration.} \Cref{lst:impl-loop} shows the experiment loop. The controller is called
every second step, which gives the 1~s decision period with a 0.5~s simulation step. Signal states are
set only through the phase manager, which inserts amber and all-red. The same loop drives B1--B3,
\prism{} and PRISM-GT, selected from \texttt{scenarios.yaml}, so the comparison is like-for-like.

\begin{lstlisting}[style=py,caption={Simulation loop (\texttt{experiments/run\_experiment.py})},label={lst:impl-loop}]
def run(scn: Scenario, seed: int, gui: bool = False) -> Path:
    sim = traci if gui else libsumo
    sim.start(["sumo-gui" if gui else "sumo", "-c", scn.cfg, "--seed", str(seed),
               "--step-length", "0.5", "--lateral-resolution", "0.8",
               "--tripinfo-output", scn.out(seed, "tripinfo.xml")])
    sensor = GroundTruth(sim) if scn.oracle else VirtualLidar(sim, seed=seed)
    tracker = PerceptionTracker(alpha=0.5, beta=0.1, gate=2.5)
    ctrl = make_controller(scn.controller, PhaseManager(sim, TLS_ID))
    log, step = DecisionLog(scn.out(seed, "decisions.csv")), 0
    while sim.simulation.getMinExpectedNumber() > 0 and step < 2 * 3900:
        sim.simulationStep()
        step += 1
        if step % 2:                      # decide every 1 s (two 0.5 s steps)
            continue
        t = sim.simulation.getTime()
        objs = tracker.update(sensor.scan(), t)
        decision = ctrl.step(objs, t)     # PRISM: decide(); B1/B3: plan; B2: SUMO
        log.write(t, decision)
        if gui:
            show_overlay(sim, objs, decision)   # VLS points and reason
    sim.close()
    return scn.out(seed, "tripinfo.xml")
\end{lstlisting}

\paragraph{SUMO configuration.} \Cref{lst:impl-xml} gives an excerpt of the vehicle types and the B1
program. The fail-safe re-uses exactly this program, so the fallback is the plan that road users
already know.

\begin{lstlisting}[style=xml,caption={Excerpt of \texttt{vtypes.add.xml} and the B1 signal program},label={lst:impl-xml}]
<vType id="motorcycle" vClass="motorcycle" length="2.0" maxSpeed="19.44"
       accel="2.8" sigma="0.5" tau="1.0" lcSublane="10" minGapLat="0.3"
       speedFactor="normc(1,0.1)"/>
<vType id="threewheeler" vClass="passenger" length="2.7" width="1.4"
       maxSpeed="13.89" accel="1.8" sigma="0.5" tau="1.0" guiShape="passenger/hatchback"/>
<vType id="bus" vClass="bus" length="11.5" maxSpeed="16.67" accel="1.2"/>
<tlLogic id="maviththara" type="static" programID="B1_fixed78" offset="0">
  <phase duration="21" state="GGrrrrGGrrrr"/>  <!-- phi1 bypass through -->
  <phase duration="3"  state="yyrrrryyrrrr"/>
  <phase duration="2"  state="rrrrrrrrrrrr"/>
  <phase duration="8"  state="rrGrrrrrGrrr"/>  <!-- phi2 bypass right   -->
  <phase duration="3"  state="rryrrrrryrrr"/>
  <phase duration="2"  state="rrrrrrrrrrrr"/>
  <phase duration="15" state="rrrGGrrrrGGr"/>  <!-- phi3 cross through  -->
  <phase duration="3"  state="rrryyrrrryyr"/>
  <phase duration="2"  state="rrrrrrrrrrrr"/>
  <phase duration="14" state="rrrrrGrrrrrG"/>  <!-- phi4 cross right    -->
  <phase duration="3"  state="rrrrryrrrrry"/>
  <phase duration="2"  state="rrrrrrrrrrrr"/>
</tlLogic>
\end{lstlisting}

% ------------------------------------------------------------------
\section{Chapter summary}

This chapter described how the artefact was designed and built. \dtx{} turned the drone videos into
anonymous metric trajectories through stabilisation, YOLOv8s detection with SAHI, ByteTrack tracking,
homography and Savitzky--Golay smoothing, and it measured waiting time with the same 0.1~m/s
definition that \sumo{} uses. The \sumo{} model reproduced the left-hand, 2+2-lane junction with
give-way slip lanes, sublane motorcycles and a four-phase signal, first with mock demand and then with drone data. The \vls{} supplied a noisy, occluded, class-blind
object list. \prism{} combined PCU-weighted movement pressure, an arrival term, a downstream penalty
and a quadratic urgency factor in six transparent rules, with proven safety, a 175~s fairness bound,
$O(n)$ cost and a logged reason for every decision. Three baselines were implemented, including a
Webster plan that re-times the signal from the measured demand, together with a ground-truth oracle.

\Cref{ch:eval} now tests whether these designs meet the requirements and how \prism{} performs
against the baselines.
